% Section 2 -- Related Work. Pages target ~1.0.
% Must cite: PaperOrchestra, Sibyl, AI Scientist v2, Karpathy autoresearch,
%            Agent Laboratory.
\section{Related Work}
\label{sec:related}

We position \texttt{cursor\_manager} along two axes: \emph{scope} (does
the system attempt end-to-end autoresearch, or only writing?) and
\emph{evaluation target} (does the system optimize for a benchmark
scoreboard, or for a real-conference submission?). To our knowledge no
prior published system occupies the ``writing-only \(\times\)
submission target'' quadrant in which we operate.

\paragraph{End-to-end autoresearch frameworks.}
\citet{lu2024aiscientist} introduced the AI Scientist as a
single-process pipeline that ideates, implements experiments, and writes
the resulting paper, evaluated on three machine-learning subfields with
held-out reviewer LLMs as the scoring oracle. The successor system
\citet{yamada2025aiscientistv2} introduces an agentic tree-search
exploration step and is the first end-to-end framework to obtain a
peer-reviewed acceptance: one of three submitted papers was accepted at
the ICLR 2025 ``I Can't Believe It's Not Better'' workshop after a
standard review process.\footnote{The remaining two submissions of
\citet{yamada2025aiscientistv2} were not accepted; the authors report
that the workshop was made aware of the AI authorship in advance, which
is a venue-cooperation pattern that we do not assume in the
\texttt{cursor\_manager} deployment.} \citet{schmidgall2025agentlaboratory}
present Agent Laboratory, a three-stage pipeline (literature review,
experimentation, report writing) optimized for cost reduction rather
than acceptance. \citet{tang2025airesearcher} extend the line with
end-to-end innovation evaluation.
\citet{karpathy2026autoresearch} releases a single-GPU framework
focused on iterating against the \texttt{nanochat} training pipeline.
\citet{sibyl2026system} implements a 19-stage pipeline on Claude Code
with $\geq$~20 specialized agents and a self-evolution loop, drawing
explicit inspiration from FARS, the AI Scientist line, and
\citet{karpathy2026autoresearch}.

\paragraph{The throughput axis: FARS.}
The closest published deployment in optimization target, and the
most important contrast for our positioning, is
\citet{fars2026analemma}. FARS is the first publicly demonstrated
industrial-scale autoresearch system: in a February 2026 livestreamed
deployment it produced 100 short papers in 228 hours of unattended
operation on a 160-GPU cluster, consuming 11.4 billion tokens at
roughly 104 thousand US dollars. The FARS-100 papers scored a mean
of 5.05 on the Stanford Agentic Reviewer calibrated to ICLR scoring,
which is 0.84 above the ICLR 2026 human submission mean of 4.21 and
0.34 below the ICLR 2026 acceptance threshold of 5.39 (the mean over
5340 of 19814 accepted submissions). FARS is a deliberate
scale-throughput design: short single-contribution papers explicitly
including negative results, with no commitment to crossing any
specific venue's acceptance threshold. cursor\_manager occupies the
orthogonal axis: rather than maximizing throughput at
borderline-reject quality, we ask what discipline mechanisms move a
single in-progress paper from the 5.05 band into the accept-decision
band of a top-tier venue's primary track. The two are complementary;
a FARS draft could in principle be routed through a
cursor\_manager-style discipline layer, but we do not evaluate that
composition.

\paragraph{Manuscript-writing-specialized agents.}
\citet{song2026paperorchestra} present PaperOrchestra, the closest
prior work in scope: five specialized agents (Outline, Plot, Literature
Review, Section Writing, Content Refinement) compose a writing-only
pipeline evaluated on PaperWritingBench, a benchmark of 200
reverse-engineered CVPR / ICLR 2025 papers. The PaperOrchestra
literature-review agent inspires the design of our
\texttt{lit\_review} sub-agent
(\S\ref{sec:discipline-rules}).
PaperOrchestra's evaluation, however, scores against the held-out
ground-truth reverse-engineered paper and does not include any
real-submission deployment, leaving open whether the benchmark-grade
pipeline transfers to genuine submission constraints. The recent
benchmark of \citet{miyai2026paperreconstruction} extends this
benchmark-only paradigm with an explicit hallucination-detection
component but again does not include real-submission deployments.

\paragraph{Discipline mechanisms in agent systems.}
The discipline mechanisms we introduce in \S\ref{sec:discipline} have
partial precedents in prior systems but none combines them in the
form we deploy. \citet{sibyl2026system} use a multi-agent debate
quality gate, which is a confidence-aggregation pattern rather than a
hard-rule pattern; debate can in principle be voted-down by the same
reward-hacking dynamic it is meant to check.
\citet{yamada2025aiscientistv2} use agentic tree search to backtrack
when an experiment fails, which is closer to our Sentinel watchdog in
spirit but operates inside the experimentation stage rather than as a
deployment-level monitor. PaperOrchestra's content-refinement loop is
quality-improving rather than claim-protective: it can rewrite text to
be smoother but does not enforce that the rewrite preserves the
original claim. The
\texttt{[BLOCKED-DECISION]} escalation protocol of
\S\ref{sec:discipline-blocked}, the cross-model adversarial
configuration of \S\ref{sec:discipline-crossmodel}, and the explicit
silent-failure monitor of \S\ref{sec:discipline-sentinel} are, to our
knowledge, novel in the autoresearch literature, and we ablate each
in \S\ref{sec:case}.
